\subsection{分次分解}

在本节中，假设环 A 配备了一个分次结构 \((\mathbf{A}_n)_{n\in \mathbf{Z}}\) ，使得 \(\mathbf{A}_n = 0\) 当 \(n < 0\) 。称一个分次 A-模 M 是下有界的，如果 \(\mathbf{M}_n = 0\) 当n足够小时；每个有限生成的分次A-模都是下有界的。

命题11. --- 若 M 是下方有界的分次 A-模（相应地，若 M 是有限生成的分次 A-模且 A 是左诺特的），则存在一个分次 A-模的正合序列，向左无界，

\[
\dots \longrightarrow \mathrm{L} _ {n} \xrightarrow {d _ {n}} \mathrm{L} _ {n - 1} \longrightarrow \dots \longrightarrow \mathrm{L} _ {1} \xrightarrow {d _ {1}} \mathrm{L} _ {0} \xrightarrow {d _ {0}} \mathrm{M} \longrightarrow 0
\]

其中 \(L_{i}\) 是分次自由的且有下界（相应地，分次自由且有限生成），并且其中 \(d_{i}\) 是次数为 0 的分次同态。

若 N 是下方有界的分次 A-模（相应地，且在诺特环 A 上有限生成），则存在下方有界（相应地，且有限生成）的分次自由 A-模 L 以及一个满射的分次同态 L → N（II，第 167 页，注记 3）。

既然如此，假设给定一个分次A-模与次数为0的分次同态的正合序列

\[
\mathrm{L} _ {n} \xrightarrow {d _ {n}} \mathrm{L} _ {n - 1} \longrightarrow \dots \longrightarrow \mathrm{L} _ {0} \xrightarrow {d _ {0}} \mathrm{M} \longrightarrow 0,
\]

其中 \(L_{i}, i = 0, \ldots, n\) ，它们是分次自由的且有下界（分别为分次自由且有限生成的）。则N = Ker \(d_{n}\) 下有界（相应地，有限生成）；因此存在一个有下界的分次自由A-模（相应地，分次自由且有限生成） \(L_{n+1}\) 以及一个次数为0的分次同态 \(d_{n+1}: L_{n+1} \to L_{n}\)，使得Im \(d_{n+1} = N\)；序列

\[
\mathrm{L} _ {n + 1} \xrightarrow {d _ {n + 1}} \mathrm{L} _ {n} \xrightarrow {d _ {n}} \mathrm{L} _ {n - 1} \longrightarrow \dots \longrightarrow \mathrm{L} _ {0} \xrightarrow {d _ {0}} \mathrm{M} \longrightarrow 0
\]

于是正合。该命题随后由前述内容对n归纳得出。

